\section{A divided slice and enclosing loops}\label{L:part7}

This section has two outputs: a constant-factor cylinder estimate for an
enclosing-loop gas, and a positive comparison between its plane version
and the mass of all boundary chords visiting a specified interior vertex.
The latter comparison, equation~\eqref{L:eq6}, turns the analytic
partition function into a local length observable.

On the periodic cylinder split a slice of \(m+n\) ports at two prescribed interfaces between columns into a first interval \(\mathcal X\) of \(m\) and its complement \(\mathcal Y\) of \(n\) (the rapidities are denoted by the same intervals). At the honeycomb point we can mark two hexagon centers at the split, exactly on the slice. Let \(\mathcal F(\mathcal X;\mathcal Y;\lambda)\) count diagrams consisting of disjoint loops, giving \(\lambda\) to each loop \emph{separating} the two marks, zero to every other loop. Separation is on the sphere obtained by capping both ends, so can just be tested by odd parity of occupied ports of that loop in \(\mathcal X\). We include the empty diagram with weight 1. We will use
\begin{equation}
 \mathcal F(1^m;1^m;2)\asymp m^{1/6}.\label{L:eq1}
\end{equation}

Here is a finite formula and proof. Give each site type \(\sigma_i=1\) in \(\mathcal X\), \(-1\) in \(\mathcal Y\). Define pair denominators and scalar functions (subscript given by the product \(\varepsilon=\sigma_i\sigma_j\)) by
\[
 D_+(x,y)=K(x,y)/(xy),\quad D_-(x,y)=xy,\qquad
 S_+(u)=u/[(u-q)(u-q^{-1})],\quad S_-(u)=(1+u)^2/u .
\]
Use pair factors \(P^\varepsilon_{ab}(z_j/z_i)\) of colors \(a,b\in\{+,0,-\}\) by the same rule as the periodic color sum but with \(S_\varepsilon\): \(P_{++}=P_{--}=1,\ P_{00}=S_\varepsilon,\ P_{+0}=P_{0-}=S_\varepsilon(u/q),\ P_{+-}=S_\varepsilon(u/q)S_\varepsilon(u/q^2)\), completed by transposed inversion. We can also use the conjugate convention \(q\mapsto q^{-1}\) here. Write \(t_a=1,0,-1\) on these colors, \(\rho=2-\sqrt2\). Take \(w_0^\sigma=1,\ w_\pm^\sigma=\sqrt\rho\,e^{\pm i\phi_\sigma}\), where \(2\cos(\phi_++\phi_-)=\lambda\). Then the raw sum is
\begin{equation}
 \mathcal H:=\frac{Q(\mathcal X,\mathcal Y)^2\,\mathcal F}{\prod_{i<j}D_{\sigma_i\sigma_j}(z_i,z_j)}
 =
 \sum_{\sum_i \sigma_i t_{a_i}=0}
 \prod_i w_{a_i}^{\sigma_i}
 \prod_{i<j}P^{\sigma_i\sigma_j}_{a_i a_j}(z_j/z_i).\label{L:eq2}
\end{equation}

Indeed one glues two periodic \(V^+\)'s, lower and reflected upper, at that slice (no eligible loop can be entirely in one of the halves). The parity test uses matchings only. Absolute convergence and rational continuation, hence polynomiality with \(Q^2\) and degree \(\le2(m+n-1)\) per variable, follow from the periodic vacuum bounds. There is separate symmetry and both pair recursions of the periodic calculation hold within a type (reduce \(\mathcal F\) with no multiplier). This uses the same swap cancellation and sliding of \(P\) or \(C\) across the gluing, with contraction into the reflected vacuum by the \(P P^t,C C^t\) exchange relation. All small added maps are confined to the one interval, so loops closed off and killed on one side in doing this had even parity, and through parities are unchanged. Extreme-site removal again holds. One more condition when \(m=n\) is to set all \(x_i=q^3 y_i\) simultaneously, pairing \(\mathcal X,\mathcal Y\) in reverse orders. Lower cap reductions peel from the \(\mathcal X\)-to-\(\mathcal Y\) interface, reflected upper cap reductions from the other interface. They give at each matched pair either two vacant caps or two cups closing a loop of odd parity, independently (take compatible half-power branches). Thus \(\mathcal F=(1+\lambda)^m\). These conditions determine \(Q^2\mathcal F\) inductively starting with size zero: within the larger group test any variable at the \(4(\max(m,n)-1)\) gap-two/three nodes, 0 and infinity (leading coefficient). Only at \(m=n>0\) is there a possible difference, divisible by \(\prod z_i\prod_{i<j,\ \sigma_i=\sigma_j}K(z_i,z_j)\) and hence a scalar multiple by degree, fixed by the cap condition.

To check \eqref{L:eq2}, use symmetry and residues in \(u=z_j/z_i\) near same-type positive gaps. At \(q^2\), labels \(+0,0-,+-\) fuse to \(+, -,0\); at \(q^3\), \(+-\) disappears. The condition on total charge is preserved with bijective reductions for each fixed choice. Spectator factors relative to the reduced ones are \(S_\varepsilon(z/(qz_i))\) and \(S_\varepsilon(z/(qz_i))S_\varepsilon(z/(q^2z_i))\) respectively, exactly those of \(Q^2/\prod D\) by the periodic table. The relative pair pole coefficients (coordinate \(u/u_0-1\)) agree also: \(1/(q-q^{-1})\) for fusion and \(\rho S_+(q^2)/(q-q^{-1})\) for gap three. At \(q\) the \(00,+-\) spectators agree for every color and type and residues cancel using \(\rho S_+(1)=1\); at 1 use same-type transposition. No cross-type poles are needed. On multiplying the sum by \(\prod D\), at \(z_i=0\) each term has order \(t_{a_i}^2\) by neutrality; at infinity the degree bound is similarly improved by \(t_{a_i}^2\). For \(a_i=0\) the removal multipliers are precisely \(\prod_{j\ne i}z_j^2\) and 1 respectively (shifts contribute total power zero of \(q\), by the same neutrality). This proves the polynomial bounds and removal rules. At paired specialization the cross zeros force \(t_{a(x_i)}\le t_{a(y_i)}\), thus equality on every pair. Distinct pairs then decouple by \(S_+(v)^2 S_-(q^3v)S_-(q^{-3}v)=1\); each pair contributes \(\rho(1+\lambda)\). Meanwhile \(Q^2/\prod D=\rho^m\) by the cap recursion (spectator factors cancel between pairs). This proves the formula. Conjugating \(q\) gives identical checks.

\subsection{Doubled change of expansion}\label{L:step10}

The constraint $\sum_i\sigma_i t_{a_i}=0$ in~\eqref{L:eq2}
introduces a fixed axis in the original color channel. After conjugation,
the root channel has a different fixed axis. The construction below
identifies both axes and the momentum sums they permit. Its convergence
proof will also differ from Section~\ref{L:step6}: short-root interactions
alone are not the object bounded there; part of the kinetic propagation
must be included. Use \(n=m\), the conjugate convention, and \((\phi_+,\phi_-)=(0,0)\) or \((\pi/4,\pi/4)\). The following data replace the one-family vectors of the scalar change of channel above. Use \(D_j^\sigma\), \(j\bmod8\), modulo the radical, rotated by \(O\), with rows
\[
 D_0^\sigma D_j^\sigma=(2,-1,1,0,0,0,1,-1),\qquad
 D_0^\sigma D_j^{-\sigma}=(0,0,0,-2,2,-2,0,0).
\]
Put \(D_0^\sigma=(1-O+O^2)a^\sigma,\ a_s=a^++a^-,\ \delta^\sigma=D_7^\sigma,\ d_*=\delta^++\delta^-\), and \(M_*=O R_{\delta^+}R_{\delta^-}\). Use \(S=(1-O)^{-1}\) with value \(1/2\) on the fixed space. The old and new fixed axes and squares are
\[
 h_o=\tfrac14\sum_j(D_j^+-D_j^-),\ H_o=4;\qquad h=-OS d_*,\ H=h^2=9/4.
\]
Color vectors in each family are \(a^\sigma,\ a^\sigma-D_0^\sigma,\ -O^3 a^\sigma=a^\sigma-D_0^\sigma-D_1^\sigma\), paired with \(h_o\) as \(\sigma(1,0,-1)\). All intermediate momenta are taken in the full lattice dual to the root span \(\sum\mathbb ZD_j^\sigma\).

Here the translation cocycle in a field \(V_\beta(s)\) (including source and root fields) uses \(e^{\pi i \mathcal B(\beta,P)}\) in place of \(e^{\pi i S\beta\cdot P}\). The lift \(\widehat O\) has also a momentum phase \(e^{\pi i\chi(P)}\) on input. Use
\[
\begin{aligned}
 v O^j a^\sigma &=\sigma(-1,0,1,2,3,4,1,-2)_j,\qquad v'=(O^{-1}v-v-h_o)/4,\\
 \mathcal B(b,P)&=S b\cdot P+[(h_o b)(vP)-(vb)(h_o P)]/4,\qquad
 \chi(P)=(h_oP)(v' P).
\end{aligned}
\]
\begingroup\emergencystretch=1em
Some useful checks: radicals are the common sum and difference alternating sum; on the sum/difference sectors the metric matrices for contractions are the circulants adding/subtracting the two displayed rows. The \(a^\sigma\) rows on \(D_j^\sigma,D_j^{-\sigma}\) are \((1,0,0,0,0,0,1,0),(0,0,0,-2,0,0,0,0)\) respectively. We have \(hD_j^\sigma=(-1,0,-1,1,-1,1,0,1)\), \(ha^\sigma=0\). The determinants on the quotient are
\[
 \det(I-yO)=(1-y^8)^2/(1-y^2),\quad
 \det(I-yM_*)=(1-y^6)^4(1-y)^2/[(1-y^3)^2(1-y^2)^3].
\]
Indeed before quotienting the common and difference \(M_*\) determinants are \((y-1)^2(y^2-y+1)(y^2+y+1)^2\), \((y+1)^2(y^2-y+1)^3\). In particular no Jordan derivative on a fixed axis is needed. \(\mathcal B+\mathcal B^t\) is the vector metric, and \(\mathcal B\) of any root or color against a root in either order is integral. For instance \(vD_j^\sigma=\sigma(0,1,2,3,0,1,2,-1)_j\); the two rows of \(\mathcal B(a^\sigma,D)\) are \((1,1,1,1,0,0,1,0),(0,0,0,-2,-1,-1,-1,0)\). These checks follow by multiplying the circulants, using \(S=-\sum jO^j/8+4\,{\rm proj}_{\rm fix}\). For \(v'\beta=0\), in particular roots of indices \(7,0,1\), \(\widehat O V_\beta\widehat O^{-1}=V_{O\beta}\).
\par\endgroup

Use the oscillators, degree and torus coordinates \(p=e^{-L/8},\ \mathfrak q=e^{-4\pi^2/L},\ s=e^{-2\pi i u/L}\) as before. Temporarily separate the sources (\(z_i=e^{u_i/8}\) near 1, vector \(a^{\sigma_i}\)). In the angular trace with \(\mathfrak q^D e^{-2\pi iD}\widehat O\) conjugate each source field by the product over types of
\[
 g_\sigma=\exp(\xi_\sigma F_1^\sigma)\exp(\xi_\sigma F_0^\sigma),
 \quad F_j^\sigma=f_{D_j^\sigma},\quad \xi_\sigma=e^{-i\phi_\sigma}.
\]
As before \(f_b=[V_{-b}]_0,\ e_b=-[V_b]_0\) use zero modes. Their residue commutators and Weyl action hold with \(S b\cdot P\) replaced by \(\mathcal B(b,P)\) (the required integral exchange signs follow from the preceding identities). On a separated source this gives the three color fields with coefficients \(1,\xi_\sigma,\xi_\sigma^2\); the two successive residue cocycles use \(\mathcal B(D_0,a)=\mathcal B(D_1,a-D_0)=0\) within that family, and the opposite family commutes here.

Equivalently remove the conjugation, replacing the background at the cut by \(\widehat M_*\) and root zero-mode exponentials distributed around the trace. Here \(\widehat M_*=\widehat O W_+W_-\) uses the same root Weyl matrices as above with parameters \(\xi_\sigma\); on momentum \(P\), each \(W_\sigma\) uses multiplier
\[
 (-\xi_\sigma)^{\delta^\sigma P} e^{-\pi i(\delta^\sigma P)\mathcal B(\delta^\sigma,P)}
\]
and reflects vectors and oscillators in that root. (The two cross cocycles are zero.) More precisely, suppress both family superscripts and multiply the indicated factors over families; write \(\widetilde B=\exp(-\xi F_0)\exp(\xi F_1)\exp(\xi F_0)\), \(B=\widehat O\widetilde B\widehat O^{-1}\), \(X_-=\exp(\xi f_\delta), X_+=\exp(-e_\delta/\xi)\). Cycling the conjugators gives background \(\widehat O\exp(-\xi f_\delta)\widetilde B\). Commute \(\widetilde B\) left through the \(f_\delta\)'s (it uses \(-D_1,-D_0-D_1\), products nonnegative), then cycle as \(B\) and commute through the sources. Now use \(\widehat O\exp(-\xi f_\delta)=\widehat M_* X_- X_+\). Move \(X_+\) across \(B\) by conjugating \(B\), then through the sources. Thus in order of action the root circles have the following vectors in each family and reduced angular coordinates (\(u=x+2\pi i\,d_i\)):
\begin{equation}
\begin{array}{c|ccccc}
 r_i&\delta&-D_2&-D_1-D_2&\delta-D_1-D_2&-\delta\\ \hline
 d_i+d_0&0&1/2&1/2&1/2&d_f\\
 h r_i&1&1&1&2&-1
\end{array}\label{L:eq3}
\end{equation}
Here take \(d_0=1/4,\ d_f<1\) sufficiently close to 1; sources lie just after the first layer. Indeed \(B\) uses two commuting roots \(-D_2,-D_1-D_2\) by the single-pole bracket; conjugating it by \(X_+\) adds only the displayed composite (product \(-1\) with the second in the same family), with no further brackets. Products within each simultaneous layer are nonnegative. All exponential coefficients are fixed constants from these group actions (their precise bracket signs are immaterial).

We describe scalar evaluations, using the short kernels \(K_U\), self constants \(d_U\), \(T_U=(1-U^{-1})^{-1}\) with value \(1/2\) on invariants, and oscillator determinants \(Z_U\) from the Gaussian change of channel above. For a list of charges \(\beta_i\) in angular order, \(b=\sum\beta_i\), take \(U=O\) (color evaluation) or \(M_*\) (root-circle evaluation), fixed vector \(h_U=h_o\) or \(h\), square \(H_U\). Require \(h_U b=0\); there can also be a parity restriction as below. Use exactly the short cross, self/image, and kinetic factors previously given with \(h_U,H_U\); in particular the kinetic exponent in \(x\)-order is
\[
 -\sum_{\rm gaps}\Delta u\, E_U(f+\sum_{\rm before} h_U\beta_i),\qquad
 E_U(y)=y^2/(2H_U).
\]
Use the prefactor \(Z_U\sqrt{L/(2\pi H_U)}\). The allowed \(f\) run over \(\mathbb Z\) or \(f_0+\mathbb Z,\ f_0=-(\phi_++\phi_-)/(2\pi)\), respectively. Phases include the same radial-order phase \(-\sum_{i<j}\operatorname{sgn}(x_j-x_i)\beta_i T_U\beta_j\) in units of \(\pi\), and \(-2 f h_U P_0/H_U\) where \(P_0\) is a representative of angular cut momenta. Other phases are the position-independent angular ones evaluated there, including translations between fields. Indeed the free Gaussian transformation only uses \(P=-T_U b\) on the perpendicular complement, the axis Gaussian summation with unit steps in \(h_U\), and the geometric oscillator traces. We give the phase and lattice checks to justify this use.

For \(O\) the cut condition \(P=O(P+b)\) gives \(P=OS b+\gamma h_o,\ \gamma=n'-vb/4\), exactly \(n'\in\mathbb Z\) by \(\mathcal B(b,D)\) integral. Then
\[
 -P^2+\mathcal B(b,P)+\chi(P+b)=-b^2/2-8\gamma n' .
\]
Thus the angular phase including field translations is \(-\sum\Delta_{\beta_i}-\sum_{i<j}\mathcal B(\beta_i,\beta_j)\) modulo 2 (\(vb\) integral). For \(M_*\), write \(n_\sigma=\delta^\sigma(P+b),\ k'=\sum n_\sigma\delta^\sigma\), \(e'=b-k'\). We must have \(P=O(P+e')\), \(h_o e'=0\), so \(n_+-n_-=h_o b\) and \(P=OS e'+\gamma h_o,\ \gamma=l-v e'/4,\ l\in\mathbb Z\). Self-consistency on the two \(\delta\)'s now says \(h b=0,\ 2l=-\mathcal B(b,\delta^+-\delta^-)\), by \(S d_*=d_*-h\). Thus either no terms occur or the \(n_\sigma\)'s step together freely with \(P\) stepping by \(h\). The phase part
\[
 -P^2+\mathcal B(b,P)-\mathcal B(k',P+b)+\chi(P+b-k')
    \equiv -{e'}^2/2-\mathcal B(k',b)\pmod 2
\]
as in the \(O\) case. With \((-1)^{\sum n_\sigma}\) included this gives \(-b^2/2+\mathcal B(b,k')\), unchanged under the step since \(\mathcal B(b,d_*)=hb\). Thus only \((\xi_+\xi_-)^{\rm step}\) modulates the summation. Ordinary Poisson transformation gives the specified indices. All residual phases, for the two parameter choices, depend on integral charge multiplicities modulo fixed integers (and on the two orders), using rational bilinear data.

Let \(\mathcal T\) denote the resulting fused trace at \(u=0\), vector \(\alpha=m a_s\) before group action. Recovery at confluence uses
\begin{equation}
 \mathcal H(1^m;1^m;\lambda)
  =[p^0]\,
 \frac{\rho^m e^{im(\phi_++\phi_-)}\mathcal T}
 {Z_O\sqrt{L/(2\pi H_o)}\ e^{-\pi i\Delta_\alpha}
 d_O(\alpha)(L/(2\pi))^{\Delta_\alpha}} .\label{L:eq4}
\end{equation}
To check, for separated sources divide by the leading uncolored radial prescription at \(O\), i.e. with \(f=0\) and no finite-period images. Choose angular and real orders identical. The ratios of \(d_O\) on the three colors are \(1,1/\sqrt\rho,1\), and dimensions agree. The relative pair phase at \(f=0\) is \((\mathcal B(\beta_i,\beta_j)-S\beta_i\cdot\beta_j)-(\mathcal B(a^{\sigma_i},a^{\sigma_j})-S a^{\sigma_i}\cdot a^{\sigma_j})\) modulo 2 by integrality of the differences. Since \(v\beta_i=\sigma_i(-1,-1,-2)_{a_i}\), this gives \(q^{-\sigma_i\sigma_j b_{a_i a_j}}\), \(b_{+0}=b_{0-}=1,\ b_{+-}=3\), completed antisymmetrically. The relative kinetic multiplier is \((z_j/z_i)^{\sigma_i\sigma_j(t_{a_i}t_{a_j}-1)}\) per pair by neutrality. Including the direct kernels with factors \((1-q^{-j'} z_i/z_j)^{\beta_j O^{j'}\beta_i}\) relative to the uncolored ones gives precisely \(P^{\sigma_i\sigma_j}\) in conjugate convention by substitution. The normalization therefore yields \eqref{L:eq2}. All other terms use image factors regular at \(p=0\) and/or \(p^{f^2},\ f\ne0\). Confluence gives \eqref{L:eq4} exactly as for the earlier torus formula: the leading local product factor of uncolored fields is common upstairs and in the leading denominator by vertex contraction (conjugation preserves it). In the transformed expression root layers stay clear of the sources. In the color ratio one can take the constant coefficient in the separation, with uniform Laurent bounds near \(p=0\) since direct ratios are meromorphic with bounded pole orders and all image series are regular. Thus that coefficient commutes with \([p^0]\).

\subsection{Convergence check}\label{L:step11}

The estimate needed for summing the doubled trace is the following.
Let $\mathscr W_{\rm root}$ denote the product of the root short-range
factors and root self constants/images in the $M_*$ prescription. Omit
the root field powers of $L/(2\pi)$, which cancel against zero-mode
measures, and omit all source contractions.

\begin{lemma}[Doubled root stability]\label{L:lem:doubled-stability}
For the ten root species and angular layers in~\eqref{L:eq3}, one can
choose $d_f<1$ sufficiently close to $1$ and constants $0<\eta<1$ and $c,C>0$
so that the following holds. If the total root $h$-charge is zero,
$K$ roots lie on a real circle of sufficiently large period $L$, and
$K_I$ counts them in successive bins of lengths between $1/2$ and $2$,
then, for every allowed axis index $f$,
\[
 |\mathscr W_{\rm root}|\,
 \exp\!\left[-(1-\eta)\sum_{\rm gaps}
       \Delta x\,E\Bigl(f+\sum_{\rm before}hr_i\Bigr)\right]
 \le \exp\!\left(CK-c\sum_I K_I^2\right),
 \qquad E(y)=\frac{y^2}{2H}.
\]
The constants are uniform for the common small-slope bends used in the
change of channel, with real gap lengths in the kinetic modulus. They
are locally uniform at bounded positive real periods as well, with bin
lengths bounded above and below by positive constants chosen for the
compact period interval.
\end{lemma}

\begin{proof}
For the ten \emph{signed as in} \eqref{L:eq3} root species let
\[
 P_{ij}(t)
 =r_i\{ e^{-a t}(1-e^{-t}M_*^{-1})^{-1}
       + e^{a t} e^{-t}M_*(1-e^{-t}M_*)^{-1}\}r_j,\qquad a=d_j-d_i\ge0,
\]
made symmetric, and write \(\mathbf v_i=h r_i\). The real short covariance has symbol \((2\pi^2/L)(P-2\mathbf v\mathbf v^{\mathsf T}/(tH))/t\), \(t=2\pi|2\pi l/L|\), as in the earlier stability argument. We claim that for some small \(\eta,c>0\)
\begin{equation}
 P-2\eta\mathbf v\mathbf v^{\mathsf T}/(tH)
      -\mathscr B\,\tanh(3t/2)\ \ge\ c\min(t,1)I\label{L:eq5}
\end{equation}
for a symmetric array \(\mathscr B\) strictly positive \emph{entry by entry}. To check, start with no same-type subtraction, opposite-type block
\[
 B_c=\begin{pmatrix}
0&1&4&4&0\\1&2&3&6&0\\4&3&4&8&0\\4&6&8&12&0\\0&0&0&0&0
\end{pmatrix}/4
\]
in \(\mathscr B\), \(\eta=0,\ d_f=1\). Block diagonalize into common and difference five-by-five arrays \(P_\pm\). The following polynomial check is included for the strict bounds. Set \(b=2\cosh(t/2),\ e=b^2,\ g=e-3,\ j=e-2,\ l=e^2-4e+2,\ l'=3e^2-9e+4\); put \(L_+=(e-4)(e-1)^3,\ L_-=(e-1)g^2\). Multiplication gives \(P_\pm/\tanh(3t/2)=G_\pm/L_\pm\). Here are upper triangles:
\[
\begin{array}{l}
G_+:\\
2 e j l,\ b g(2e^2-5e+4),\ b(2e^3-12e^2+23e-12),\ b g j(5e-4),\
 -e(e^3-5e^2+7e-4)\\
2e^2g^2,\ e^2g j,\ e g l',\ -b g l'\\
2e j l,\ e g l',\ -b^3 j^2\\
2e g l',\ -b g(4e^2-9e+4)\\
2e j l\\[2pt]
G_-:\\
2j l,\ -b g(2e-7),\ -b(6e^2-36e+53),\ -b g(3e-10),\ -e^3+5e^2-3e-8\\
2(e-4)g^2,\ (e-4)g j,\ -e g^2,\ b g^2\\
2j l,\ -e g^2,\ b j(3e-10)\\
-2e g^2,\ b g\\
2j l
\end{array}
\]
For example to get these arrays use coefficient columns \(r_i\) in a single cyclic basis, metric \(C\) the sum/difference circulant, and in the two inverses use the cyclic shift times \(1-\delta\delta^{\mathsf T}C\). The polynomial matrices \(4(G_\pm\mp L_\pm B_c)\) in \(b-2\) have positive definite coefficients except the constant common coefficient \(128\,v_c v_c^{\mathsf T}\), \(v_c=(1,1,1,2,-1)^{\mathsf T}\). For clarity here are leading-minor checks by substitution (entries listed as lower endpoints over the coefficient columns):
\[
\begin{array}{c|c|rrrrr}
 &\text{powers}&1&2&3&4&5\\\hline
 +&1,\ldots,8&8&47&92&336&1344\\
 -&0,\ldots,6&8&79&452&720&2752
\end{array}
\]
These numbers can be obtained by forming for power \(a\) the entries \(\sum_{i=a}^{8}\binom ia 2^{i-a} [b^i](4G_{\pm,st}\mp (4 B_{c,st})L_\pm)\) from the upper triangles and taking successive top-left determinants (the constant common one is by direct evaluation).

The coefficient minors prove positivity for every $b>2$.
Near $t=0$, $b-2$ has order $t^2$ and $\tanh(3t/2)$ has order $t$.
The common constant coefficient is supported on $v_c$; on its
complement the first positive coefficient therefore gives order-$t$
positivity after division by $L_+$. The difference denominator stays
nonzero there, giving the same order from its positive constant
coefficient. This proves the required strict bounds at zero as well as
at positive frequencies. The expansions there, after the possible \(1/t\) term independent of angle, are \(O(t)\) uniformly with smoothly varying coefficients for \(d_f\) nearby (also by oddness of the displayed resolvents combination). Thus one can take \(d_f<1\) sufficiently close and small \(\eta>0\). At infinity the only limiting entry affected by this near-endpoint perturbation in each block is the \(1,5\) entry (and transpose), interpolating from its value to 0; all other off-diagonals in that last row are zero, so both extremes are definite, hence uniformly also the interpolates. Finally add a sufficiently small strictly positive coefficient array (e.g. all entries equal) to the original \(\mathscr B\).

Here is how we use \eqref{L:eq5}. Keep \(\eta\) of the kinetic propagation for separate estimates, borrowing \(1-\eta\) for stability. For neutral total root axis charges (the sources separately have zero \(h\)-charge), the nonconstant part of the kinetic quadratic form supplies Fourier covariance \(1/(LH(2\pi l/L)^2)\) times \(\mathbf v\mathbf v^{\mathsf T}\) before scaling by \(1-\eta\), by integrating the square of the cumulative charges. At frequency 0 use the symbol only on neutral totals. The residual real form thus dominates a positive multiple of \(1/(L(1+t^2))\) diagonally after the subtraction in \eqref{L:eq5}. The subtracted kernel has nonnegative spatial entries: \(\tanh(3t/2)/t\) is a positive sum of exponential-covariance symbols by the elementary half-integer cotangent series for \(\tanh x/x\). Near zero in space these entries are \(b_{ij}'\log(1/|x|)+O(1)\), all \(b_{ij}'>0\), by the high-frequency asymptotic, uniformly for large real period by periodization. The actual short kernels have analogous coefficients \(b_{ij}''\ge0\) (zero or \(r_i r_j\) at equal angles), with bounded remainders and derivative bounded by \(C/|x|\); away from 0 they decay exponentially by circular image distance. This uses the same \(K_U\) series and its Jordan derivatives, with only simultaneous-layer contacts.

Smooth the residual by averaging point charges in a tiny interval of size \(s_*\) (fixed). Its full form still dominates a bin-square penalty \(c'\sum_I K_I^2\) by the corresponding smoothed exponential spatial covariance; \(K_I\) counts particles in unit-size bins. Kinetic variance can cover its smoothed version including diagonals by Parseval. Diagonals of the smoothed short and subtracted kernels cost \(O_{s_*}(1)\) per particle. Off diagonal, the unsmoothed actual short kernel bounds the smoothed difference entrywise for distances \(\le\sqrt{s_*}\), since the difference equals \((b_{ij}''-b_{ij}')\log(1/(|x|\vee s_*))+O(1)\), whereas the actual singular coefficient is nonnegative. At greater distances the loss is at worst \(O(\sqrt{s_*})\) with exponential tails (discard the subtracted positive entries). Thus it is absorbed by the bin penalty. Altogether, including self constants/images of roots, the modulus excluding source contractions and with the reserved \(\eta\)-kinetic factor taken out is at most
\[
 \exp(CK-c'\sum_I K_I^2)
\]
at \(K\) particles (field power prefactors for zero modes to be canceled against measures). This works locally also at bounded positive period. As before we can vary \(\Im L\) along contours with common small-slope bend, since the change of actual short logs costs a small multiple of the bin estimate by kernel differentiation, and real gap lengths govern kinetic modulus.

\end{proof}

\paragraph{Finite coefficient support and analytic continuation.}
We now apply the estimate to justify the formal conjugations and the
summation of the trace. These are separate from the finite covariance
calculation above.

For the formal trace step first fix cut degree and source mode degrees. Before the Weyl change, on cycling \(g^{-1}\) past the cut, all root shifts in the closure with \(O\) are negative counts from indices \(7,0,1\) of the two families. The cut then has proper support: on these six columns the Gram of \(OS\) projected off \(h_o\) is strictly copositive (diagonal \(11/16\); same-family off-diagonals \(-3/16\) at distance one and \(3/16\) at distance two, opposite-family off-diagonals bounded below by \(-1/16\)). Thus degree bounds the counts, fixed-axis momentum and oscillator degree there. Cycling \(B\) after commuting as specified just rotates finite matrix paths of this kind. After that cycle the Bruhat and residue-commutation identities need only be used per fixed input vector, locally finite. The final expanded prescription is itself proper: its \emph{time-average} angular momentum norm controls a strictly positive quadratic in the counts up to \(O_m(1)\), and also the square along the freely summable axis (around a count-dependent center). Indeed with homogeneous root data the solution of the jump profile having axis mean zero is the \(t\downarrow0\) limit of solving \((\partial_d+t)y=\sum r_i\delta_{d_i}\) with twist taking the end to the start by \(M_*\). Integration by parts, or the two directed exponential propagators, gives \(\int y^2\) as the count contraction against \(P(t)/(2t)\). Thus \eqref{L:eq5}, with nonnegative entries added back, gives the assertion (external charges only translate by bounded profiles for fixed \(m\)). Angular degree between positions differs from cut degree by fixed source modes only. So along any formal coefficient path the time-average oscillator degree is nonnegative, which bounds the time-average momentum norm from above at fixed cut and source modes. Counts, free axis momentum and cut oscillator degree are therefore bounded there. This proves final support finiteness (together with the earlier cycling bound it suffices: once all cycling is done, equality of the remaining operators on each fixed input vector at fixed modes can be used before summing over that vector); zero-mode exponentials and Weyl operations per vector used here are the locally nilpotent/sliding identities of the torus lemma. It also proves absolute convergence after Wick contraction for small \(\mathfrak q>0\): at radii prescribed by the angles the momentum power of \(\mathfrak q\) is exactly half that time average, and oscillator Taylor factors cost at worst \(\exp(C_m(1+K^2))\) absolutely (direct factors within simultaneous layers are polynomials). Thus coefficients can be grouped to implement contour zero modes, using separated sources first.

After Gaussian change the bin estimate gives absolute summability for every positive real period, with Gaussian tails in \(|f|\gg 1+K\); fixed external fields off the root circles pose no problem. We can therefore continue equality along this axis and to the bounded imaginary period shifts at large real period, as in the torus argument, then fuse sources by the local product rule. More explicitly, for continuation from a given real period parameterize integration positions by their real coordinates plus \((L-\Re L_{\rm initial})\) times a fixed real bend profile, plus \(2\pi i d_i\). Use a profile zero at the external positions and Lipschitz with increments of one per period (small derivative at large real period, or just take small enough increments of the parameter locally at bounded period to ensure small slopes after bending). Orders of real coordinates are preserved in a neighborhood in real part and convergence is locally uniform along the required imaginary changes, so the piecewise kernel integrals continue analytically. Simultaneous layer contacts have only nonnegative direct powers; no poles are moved through external fields. These arguments justify equality used in \eqref{L:eq4}.

\subsection{Leading projection}\label{L:step12}

Cauchy average \eqref{L:eq4} with \(\Re L=2\ell,\ \ell=6\log m,\ |\Im L|\le8\pi\). The fused source self ratio of \(d_{M_*}\) to \(d_O\) is \(\mathcal R^{2m^2}\), \(\mathcal R=(2k/3)^2/\sqrt{2k}\), since the source is common (same contractions as in the earlier unprobed calculation, doubled). Moreover \(a_s(M_*^2-M_*+I)=0\) as a contraction row. The determinant and Poisson prefactor ratio is \(e^{L/32}(1+O(e^{-\ell/4}))\): in angular product notation the oscillator quotient alone is \(F_8^2F_3^2F_2^2/(F_6^4F_1^2)\), \(F_j=\prod_{n'\ge1}(1-\mathfrak q^{j n'})\), to which we apply the elementary product transformation above; the leading constant \(3/4\) cancels against \(\sqrt{H_o/H}\).

Put \(E(y)=y^2/(2H),\ E_0=E(f_0)\le1/72\). We claim that after division by
\[
 \rho^m\mathcal R^{2m^2} m^{3/8-12E_0}
\]
the result has the form \(c_{\mathfrak c}+O(e^{-\ell/24})\) with fixed constants on finitely many size congruence classes. We give truncation details. Work first before Cauchy averaging. Use real coordinates \(x\in[-\ell,\ell]\) and a common bend as above, zero except in fixed large seam widths, taking values \(\pm\Im L/2\) at \(x=\pm\ell\), with shape in seam coordinates independent of \(\ell\). Let \(s'=\ell-|x|\) denote distance from the seam.

The source profiles on the five root circles are multiples of \((0,-1,-1,0,1,1)\) for \(a_s M_*^j r_i,\ j=0,\ldots,5\); multipliers are respectively \(1,0,-1,0,-1\). Nonzero profiles thus contribute in modulus at most \(C\exp(-c m e^{-|x|/6})=C\exp(-c\exp(s'/6))\) including images. Indeed off the bend the direct product log with one unit of profile uses exponents \((0,-1,-1,0,1,1)\) on \(\left|1-e^{(-|x|+2\pi i(d_i-j))/6}\right|\). Comparing numerator and denominator gives the strict bound for unit profile at \(d_i<0\) or opposite profile at \(d_i>0\), since \(\cos(\pi(d_i-1)/3)+\cos(\pi(d_i-2)/3)\) has the sign of \(d_i\) here. In bounded seam widths and for the more distant images costs are bounded per particle. All remaining (uncoupled) particles have strictly positive integral \(h\)-charge. In particular \(K\) is at most twice the number of negative particles by neutrality, and all negatives are localized. The source self images cost \(O(1)\) in the log.

We can discard \(K>\ell^{2/3}\) and configurations with any coupled particle at \(s'>\ell^{1/4}\) with arbitrarily high exponential accuracy in \(\ell\). Indeed the bin estimate and negative-particle localization supply \(\exp(-c K^{2-1/4})\) (either a negative lies beyond distance \(K^{1/4}\), or use the bin squares up to there), retaining a fixed fraction of localization. Volumes and multiplicities cost at worst \((C\ell)^K\), as zero-mode self powers cancel length denominators in the measures. The \(f\)-sum costs \(O(1+K)\) with Gaussian tails by the reserved kinetic bound, and the divided-out empty propagation with energy \(E_0\) costs only \(e^{2\ell E_0}\).

For the retained counts the full kinetic rate on macroscopic \emph{empty} gaps can be used, not just the small reserved fraction. Shorten all cyclic gaps between sorted roots longer than \(s_0=C_1\log\ell\) to \(s_0\), for estimating the flat real root modulus. The resulting period is between \(s_0\) and \(O((1+K)s_0)=o(\ell)\) (shorten the empty circle too). Actual short-root logs and self images differ from their shortened flat values by \(o(1)\) for sufficiently large \(C_1\), before the bend: any modified image path crosses a long gap in both versions so has exponential decay, even with Jordan factors. The bend error is absorbed against the shortened bin penalty, since near pairs remain near after shortening. Apply \eqref{L:eq5} there, keeping the order and the value of \(f\) across the cut. Thus any interval of length \(R'\) with kinetic momentum different from \(f_0\) costs
\[
 \exp(-(R'-o(\ell))/9)
\]
up to subexponential factors: on the deleted portions we have the actual propagation with \[\min_{z\in(f_0+\mathbb Z)\setminus\{f_0\}} (E(z)-E_0)\ge1/9.\] The baseline \(E_0\) on the shortened part costs only \(e^{o(\ell)}\). This bound includes summation/integration since \(K\le\ell^{2/3}\), with the reserved shortened kinetic again available at \(|f|\gg K+1\). Outside the localized widths the momentum increases by positive units only. Consequently any root at \(s'\ge\ell/2\) forces such an interval of length \(\ell/2-o(\ell)\) towards one of the two sides. Likewise if the entire middle is empty but its momentum is not \(f_0\). These cases are therefore errors smaller than \(e^{-\ell/24}\).

It remains to put the retained particles in infinite seam coordinates \(\hat x=x\pm\ell\) about 0 (plus on \(x<0\)), complex positions \(\hat u=u\pm L/2\). Shapes of these paths are fixed, with \(|\hat x|=s'\). Middle momentum is exactly \(f_0\), determining \(f\) by seam charges. After removal of \(e^{-E_0 L}\), the kinetic expression now uses \(E(y)-E_0\) on this line, \(y=f_0\) at both infinities, no period dependence. Keep the direct short-root interactions in these coordinates and drop the rest (self images also); the latter paths have real lengths at least \(\ell\), so log errors \(K^2 O(e^{-(1/6-o(1))\ell})\) since minimum positive frequency for \(M_*\) is \(1/6\). For source factors keep just the first harmonic (\(1/6\)) at \(L/2\pm\hat u\) for a root, and at \(L\) for self images. These give fixed exponents since \(m e^{-L/12}=e^{-i\Im L/12}\). All other harmonics/images in these logs are exponentially smaller (error \(O(e^{-\ell/4})\) suffices), using \(|\hat x|\le\ell^{1/4}\) for coupled particles and no Jordan terms on source contractions. The retained exponents themselves obey the same localization. Thus replacement by these fixed prescriptions costs an admissible error (absolute normalized volume bounds \(e^{o(\ell)}\) here follow by the shortening argument). All remaining constants/phases depend only on seam data, \(\Im L\) and possibly the size class, including order signs via the two sides of the seam.

These fixed integrals can now be extended to the whole line and all counts with the same error. To see absolute tail bounds directly, treat dyadic blocks of furthest distance \(R\) to \(2R\) at large \(R\), estimating the direct kernels by stability on an auxiliary period, say \(8R\), with \(f_0\) outside the particles. Extra distant paths cost \(K^2 O(e^{-cR})\), absorbed by bins. The fixed bend is still allowed and any differing imaginary gap phases only use real energies in modulus. Large counts \(K>R^{2/3}\) or coupled particles beyond \(R^{1/4}\) again cost arbitrarily high exponential order. At smaller counts repeat shortening: a furthest positive particle forces an interval of incorrect momentum of length at least \(R-o(R)\) since negatives are all within \(R^{1/4}\) and the two asymptotic momenta are \(f_0\). Hence such blocks cost \(e^{-R/9+o(R)}\), summable starting at \(R=\ell/2\) with the required accuracy. The other extensions with furthest distance below \(\ell/2\) use the auxiliary period on scale \(\ell\) for the count and localization bounds. This proves convergent constant amplitudes with exponential accuracy, uniformly in the imaginary shifts, and establishes the claim by Cauchy averaging.

For \(\lambda=0\), \(\mathcal H=Q_{2m}^2/(2k)^{m(m-1)}\). The positive Fourier norm comparison in the calibration gives \(Q_{2m}\ge (2k/3)^{(2m)^2/2} k^{-m}\) directly by products, hence \(\mathcal H\ge \rho^m\mathcal R^{2m^2}\). The predicted leading power here is \(5/24<6/24\), so none of its leading constants can vanish. For \(\lambda=2\) the result for \(\mathcal H\) is at least as large as for \(\lambda=0\) by positivity of the gas; the power is \(3/8\), a difference \(1/6<6/24\). Again leading constants cannot vanish. Dividing now gives \eqref{L:eq1}.

\subsection{Plane version and flux reversal}\label{L:step13}

Write \(\mathcal Z_R(o)\) for the weight of sets of disjoint honeycomb polygons within distance \(R\) of \(o\), each strictly enclosing \(o\), with fugacity 2, including the empty set. For face centers, \eqref{L:eq1} gives the upper bound \(C R^{1/12}\) by fitting two independent such balls at the two marks of a cylinder of even circumference comparable to \(R\). Conversely, winding separators in \eqref{L:eq1} cost at most a bounded gas multiplier even discarding all avoidance: they must intersect the slice and their single-loop sum is bounded by \(\sum (h(P)+C)X^{|P|}\) modulo row translations, \(O(1)\) by the pressure estimate. Contractible separators lift uniquely around a prescribed representative of the puncture on their disk side. Those violating lift diameter \(N^{1+2\epsilon}\) at circumference \(N=2m\) cost superpolynomially small single-loop mass. Indeed for heights \(>N^{1+\epsilon}\) this is the cylinder height tail (count slice placements using the height). At smaller heights project the same lifted polygons to a cylinder in another lattice direction of period \(\lceil N^{1+3\epsilon/2}\rceil\); no collision is possible by the height restriction, while their transverse height now exceeds this new period by a power. The same tail bound applies there, since recovery of anchored lifts from loops modulo row shifts costs at most polynomially in length and period (the fixed enclosed point lies within a length of the polygon). Multiplying this small mass by twice the full gas partition bounds the loss with any such bad loop marked. The remaining contractibles lift into two independent plane gas spaces of radius \(C N^{1+2\epsilon}\) as an upper bound, by splitting according to enclosed puncture. Taking \(\epsilon\downarrow0\) gives
\(\mathcal Z_R(o)=R^{1/12+o(1)}\).
This also holds for \(o\) a honeycomb vertex (triangle center): the segment to a neighboring hexagon center avoids polygon crossings except at the vertex; at most one loop through the vertex can occur in the face gas, with bounded total weight by the plane rooted estimate. Discarding that loop thus costs only a bounded factor, besides bounded changes in radii.

There is a constant-factor comparison for a finite convex domain of triangles with a triangle at center \(v\) strictly inside (all neighbors in the domain). Let \(\mathcal Z_D(v)\) restrict the enclosing gas to this domain, and let chordal arcs in \(D\) have both ends free (sum over ordered boundary midports, weights \(X\) per center). Then
\begin{equation}
  \sum_{\omega\ {\rm visiting}\ v} X^{|\omega|}\ \asymp\ \mathcal Z_D(v)\label{L:eq6}
\end{equation}
with uniform constants. Here are details. Put \(\tau=3/8\); the complex extension weights \(X e^{\pm i\tau\pi/3}\) sum to 1. Start at an outer boundary port, keep extending from mid-edge to mid-edge unless at the boundary (after nonempty travel) or advancing toward an already visited center. The terminal collision at such a center decomposes into a prefix up to just before that visit and a detour around a simple polygon returning towards the center; reverse the detour as alternative. The arrival leg on the earlier visit is outside the polygon, and the two added turns including that visit are \(\pm4\pi/3\), so these dead contributions cancel. Summing terminal real flux gives exactly 1 from all boundary exits, using \(\cos(\tau\,{\rm turn})\); in the full convex domain this cosine is positive and uniformly bounded below by closing the simple arc along the boundary as before.

Remove now the triangle at \(v\), exposing three inner ports. From the same outer sources the identity still holds, with extra exits at these ports. Subtracting and summing over the outer sources gives
\[
 T=\sum_{\omega\ {\rm visiting}\ v}
        X^{|\omega|}\cos(\tau\,\operatorname{turn}(\omega)),
\]
where $T$ is the total real outer-to-inner flux in the punctured domain.
The sum on the right is over ordered boundary chords of the full domain.
Its cosine factors lie between $\cos(3\pi/8)$ and $1$. Reversing the
outer-to-inner paths preserves their real weights, so $T$ is also the
inner-to-outer flux summed over the three inner sources. For an inner source the extension identity works with the entire enclosing gas in the punctured domain included (path and gas disjoint). Indeed a collision detour reached from outside still cancels; reached from inside it encloses \(v\) and its added turns are \(\pm8\pi/3\), giving factor \(-2\) after summing orientations. To see the turns, the positive circuit omits at closure the bend at that center (\(+\pi/3\) from the convex side with unused leg outside, \(-\pi/3\) with unused leg inside), substituting its opposite at first visit instead. In the inside case this dead term now cancels exactly against the prefix colliding with a separate selected polygon of fugacity 2 there; these are all collisions with gas, and other gas loops remain unchanged. Thus total real exit flux from each such source is \(\mathcal Z_D(v)\). Paths reaching the outer side prohibit any enclosing loop, so contribute exactly the needed cross sum. Inner-to-inner exits have strictly negative cosine, since completing with the missing center gives a polygon with total turn \(\pm2\pi\), adding only \(\pm\pi/3\). Let $U\ge0$ be the absolute value of this inner-to-inner flux,
summed over the three sources and including the disjoint gas. The three
extension identities therefore give
\[
                         T-U=3\mathcal Z_D(v).
\]
Completing an inner-to-inner arc through the missing center multiplies
its weight by $X$ and gives a polygon through $v$. Each such polygon
has two orientations. Dropping its avoidance with the remaining gas
therefore bounds
\[
 U\le 2X^{-1}\mathcal Z_D(v)
                   \sum_{P\ni v}X^{|P|}
   \le C\mathcal Z_D(v),
\]
by the finite rooted plane polygon mass. Hence
$3\mathcal Z_D(v)\le T\le(3+C)\mathcal Z_D(v)$. Comparing $T$ with the
positive weighted chord sum proves~\eqref{L:eq6}.

